\subsection{Local moment consistency}\label{sec:limits-moments}

For a continuous-box quadratic $F(x)=\frac12x^{\T}Hx+b^{\T}x+c$
with $H\succeq-\nu I$, a bound in which large positive curvature is harmless
would replace $L$ by $\nu$. Exact positive energy suffices when all
coordinates are conditioned on one global cell. Matching finitely many
moments across separators does not.

\begin{lemma}[Globally conditioned moments]\label{lim:lem:mixture}
Let $(m,Y)$ satisfy $m\in C=\prod_i[l_i',u_i']\subseteq X$,
$\Sigma=Y-mm^{\T}\succeq0$ and $Y_{ii}\le(l_i'+u_i')m_i-l_i'u_i'$.
Then
\[
\tfrac12\langle H,Y\rangle+b^{\T}m+c\ge\OPT-\frac\nu8
\sum_i(u_i'-l_i')^2 .
\]
The same bound, averaged with the corresponding weights, holds for any
convex combination of such pairs indexed by global cells.
\end{lemma}

\begin{proof}
The left side equals
$F(m)+\frac12\langle H,\Sigma\rangle\ge F(m)-\frac\nu2\operatorname{tr}\Sigma$.
Also $\Sigma_{ii}=Y_{ii}-m_i^2\le(u_i'-m_i)(m_i-l_i')\le(u_i'-l_i')^2/4$, and
$F(m)\ge\OPT$. Averaging preserves the inequality.
\end{proof}

A global cell index can have $K^n$ values. Sparse relaxations instead keep
exact local information and match separator moments.

\begin{definition}[Local moment relaxation of order $k$]
\label{lim:def:moments}
Fix a tree decomposition and a finite partition of every $X_i$ into
intervals (cells). A feasible point consists of probability measures
$\mu_t$ on $X_{B_t}=\prod_{i\in B_t}X_i$, one per bag, such that for every
tree edge $\{t,u\}$ with separator $S$, every product $C$ of separator cells,
and every monomial $x_S^\alpha$ with $|\alpha|\le k$,
\[
\int_{\{x_S\in C\}}x_S^\alpha\,d\mu_t=\int_{\{x_S\in C\}}x_S^\alpha\,d\mu_u .
\]
Its value is $\sum_t\int q_t\,d\mu_t$, where $q_t$ is the sum of factors
assigned to $t$. Point masses at any $x\in X$ are feasible, so the optimal
value $R_k$ satisfies $R_k\le\OPT$. The occupied width $w_i$ is the largest
length of a cell of coordinate $i$ that carries positive mass.
\end{definition}

\begin{proposition}[Every finite order fails]\label{lim:prop:moments}
For each integer $r\ge1$ and rational $0<h\le1/(4r)$ there is a continuous
box QP $F_{r,h}$ with a supplied tree decomposition of maximum bag size
three, with the following properties. It has a unique minimizer and
$\OPT=1/4$. The bound $\nu\le2$ holds, and
$g_r=1/(8r(1+8(2r-1)^2))$ is a valid growth constant; both are independent
of $h$. For every choice of cells in which $[0,rh]$ lies in one cell of
width at most $2rh$ for each continuous coordinate, and $0$ and $1$ lie in
cells of width at most $\omega$ for each control,
\[
R_{2r-1}\le\frac14-\frac{2r+1}{16r},\qquad
\sum_iw_i^2\le(4r-2)\max\{2rh,\omega\}^2 .
\]
Consequently there is no function $C$ such that
$\OPT-R_k\le C(p,\bar\nu/g,k)\,\bar\nu\sum_iw_i^2$ for all instances, all
valid bounds $\bar\nu\ge\nu$ and growth constants $g$, all $k$, and all
cells.
\end{proposition}

\begin{proof}
\emph{Construction.} Use continuous variables $s$, $y_1,\ldots,y_{r-1}$,
$z_1,\ldots,z_{r-1}$ in $[-1,1]$ and controls $u_1,\ldots,u_r$,
$v_1,\ldots,v_{r-1}$ in $[0,1]$. Put $y_0=z_0=0$, $y_r=s$, and
\[
e_i=\frac{y_i-y_{i-1}}h-u_i\ (1\le i\le r),\quad
f_j=\frac{z_j-z_{j-1}}h-v_j\ (1\le j<r),\quad
f_r=\frac{s-z_{r-1}}h-\frac12,
\]
\[
F_{r,h}=2r\Bigl(\sum_ie_i^2+\sum_jf_j^2\Bigr)+\sum_iu_i(1-u_i)
+\sum_jv_j(1-v_j)+\tau(U+V),
\]
where $U=\sum_iu_i$, $V=\sum_jv_j$ and $\tau=1/(8r)$. The bags
$L_i=\{y_{i-1},y_i,u_i\}$ ($i\le r$), $R_j=\{z_{j-1},z_j,v_j\}$ ($j<r$) and
$R_r=\{z_{r-1},s\}$, with constants omitted, are arranged in the path
$L_1,\ldots,L_r,R_r,R_{r-1},\ldots,R_1$. The separator between $L_r$ and
$R_r$ is $\{s\}$, and all bags have at most three elements.

\emph{Reduction.} List the rescaled path nodes $N_0=0$, $N_k=y_k/h$
($1\le k<r$), $N_r=s/h$, $N_{r+l}=z_{r-l}/h$ ($1\le l<r$), $N_{2r}=0$. Up to
sign, the $2r$ residuals are $N_k-N_{k-1}-c_k$ with
$c=(u_1,\ldots,u_r,-\frac12,-v_{r-1},\ldots,-v_1)$; note
$\sum_kc_k=U-V-\frac12=:E$. Let $N^0(u,v)$ be the node values with every
residual equal to $-E/(2r)$, and $\delta=N-N^0$, with $\delta_0=\delta_{2r}=0$.
Since the residuals sum to $-E$, expanding the squares gives
\begin{equation}\label{lim:eq:momsplit}
F_{r,h}=\Psi(u,v)+2r\sum_{k=1}^{2r}(\delta_k-\delta_{k-1})^2,\quad
\Psi=E^2+\sum_iu_i(1-u_i)+\sum_jv_j(1-v_j)+\tau(U+V).
\end{equation}
Since $|E|<r$ and $|\sum_{l\le k}c_l|\le r$, every node of $N^0$ lies in
$[-2r,2r]$. Hence the corresponding variables $hN^0$ lie in
$[-\tfrac12,\tfrac12]$ and are feasible.

\emph{The reduced function.} With $e_2$ the second elementary symmetric
polynomial,
\begin{equation}\label{lim:eq:pi}
\Psi-\frac14=2\Pi+\tau(U+V),\qquad\Pi=e_2(u)+e_2(v)-UV+V .
\end{equation}
This follows from $E^2=(U-V)^2-(U-V)+\frac14$, $\sum_iu_i(1-u_i)=U-\sum_iu_i^2$,
and $U^2-\sum_iu_i^2=2e_2(u)$. The polynomial $\Pi$ is affine in each
variable, so its minimum over the box is attained at a vertex. At a vertex
with $U=a$ and $V=b$,
$\Pi=\frac12[(a-b)^2-(a-b)]=\frac12(a-b)(a-b-1)\ge0$. Therefore $\Psi\ge
\frac14+\tau(U+V)$, and \eqref{lim:eq:momsplit} shows that the unique
minimizer is $u=v=0$, $N=N^0(0,0)$, with $\OPT=\frac14$.

\emph{Growth and curvature.} By \eqref{lim:eq:momsplit}, \eqref{lim:eq:pi},
$U+V\ge\norm{(u,v)}^2$, and the path inequality \eqref{lim:eq:path} for
$2r$ edges,
\[
F_{r,h}-\tfrac14\ge\tau\norm{(u,v)}^2+\frac2{2r-1}\norm\delta^2 .
\]
Every node of $N^0$ is a partial sum of the $c_l$ plus a multiple $k/(2r)$
of $-E$, so it moves by at most $2\norm{(u,v)}_1$ when the controls move.
Hence $\norm{N^0(u,v)-N^0(0,0)}^2\le4(2r-1)^2\norm{(u,v)}^2$. The continuous
variables are $h$ times the nodes and $h\le1$, so the squared distance to the
minimizer is at most $2\norm\delta^2+(1+8(2r-1)^2)\norm{(u,v)}^2$. This gives
the growth constant $g_r=\tau/(1+8(2r-1)^2)\le1/(2r-1)$. The Hessian is $4r$
times a sum of rank-one Gram terms minus $2$ on the controls, so $\nu\le2$.

\emph{The relaxation witness.} On the left subtree take the law under which
$U=i$ with probability $\pi_i=\binom{2r}{2i}2^{1-2r}$ ($0\le i\le r$),
$u=(1^i,0^{r-i})$, $y_k=h\sum_{l\le k}u_l$ and $s=hi$. On the right subtree
take $V=j$ with probability $\rho_j=\binom{2r}{2j+1}2^{1-2r}$ ($0\le j<r$),
$v=(1^j,0^{r-1-j})$, $z_k=h\sum_{l\le k}v_l$ and $s=h(j+\frac12)$. Both are
probability laws, all residuals vanish on their supports, and all controls are
binary. Let $\mu_t$ be the corresponding bag marginals. Inside each subtree
the marginals come from one joint law, so they agree on every separator in
all moments. At the separator $\{s\}$ the laws of $2s/h$ are those of a
$\mathrm{Bin}(2r,\frac12)$ variable conditioned on being even, respectively
odd. Their $j$-th moments differ by a multiple of
$\sum_{k=0}^{2r}(-1)^k\binom{2r}kk^j$, which is zero for $j<2r$ as a finite
difference of order $2r$ of a polynomial of degree $j$. Thus the separator
moments of $s$ agree through order $2r-1$.

All continuous atoms lie in $[0,rh]$, and all control atoms lie in
$\{0,1\}$. With cells as in the statement, each separator variable carries
all its mass in a single cell, so the cellwise conditions reduce to the
moment equalities just proved. Hence the witness is feasible for $R_{2r-1}$.
By the symmetry $k\mapsto2r-k$, $\E U=r/2$ and $\E V=(r-1)/2$. The witness
value is therefore $\tau(r-\frac12)=\frac14-\frac{2r+1}{16r}$. There are
$2r-1$ continuous coordinates and $2r-1$ controls, which gives the width
bound.

\emph{Consequence.} Given $k$, choose $r$ with $2r-1\ge k$, so that
$R_k\le R_{2r-1}$, and fix $\bar\nu=2$ and $g=g_r$. The gap $\OPT-R_k$ is at
least $1/8$, whereas $\sum_iw_i^2\to0$ as $h,\omega\to0$.
\end{proof}

\begin{remark}[What the proposition does and does not say]
\label{lim:rem:moments}
The coordinate curvature of $F_{r,h}$ is $L=8r/h^2$, attained at the continuous variables. The
gap is far below $L\sum_iw_i^2/8$, so the example does not contradict an
error bound in terms of $L$.

The example does not survive exact matching of separator laws. If finitely
supported bag measures have equal separator marginals on every tree edge,
then sampling the root bag and then each child conditionally on its separator
defines a probability measure on $X$ with these bag marginals. Running
intersection makes the newly sampled coordinates distinct. The relaxation
value is then $\E F\ge\OPT$. The obstruction is therefore specific to finite
moment orders.

The case $r=2$ admits a four-variable form. On $[0,1]^4$ with
$0<h\le\frac12$, let
\[
F_h=8\Bigl(\frac sh-\frac{u_1+u_2}2\Bigr)^2+8\Bigl(\frac sh-\frac14-\frac v2
\Bigr)^2+\sum_{w\in\{u_1,u_2,v\}}w(1-w)+\frac{u_1+u_2+v}{16},
\]
with bags $\{s,u_1,u_2\}$ and $\{s,v\}$. Direct expansion with
$\delta=s/h-\frac18-\frac{u_1+u_2+v}4$ gives
$F_h-\frac14=16\delta^2+2[(1-v)u_1u_2+v(1-u_1)(1-u_2)]+(u_1+u_2+v)/16$.
Hence the minimizer $(h/8,0,0,0)$ is unique and $g=1/22$ is valid, using
$\norm{x-x^*}^2\le2\delta^2+\frac{11}8\norm{(u_1,u_2,v)}^2$ and
$\sum w\ge\norm w^2$. Moreover $\nabla^2F_h+2I=16r_1r_1^{\T}+
16r_2r_2^{\T}+2e_se_s^{\T}$ with $r_1=(1/h,-\frac12,-\frac12,0)$
and $r_2=(1/h,0,0,-\frac12)$. Since $(0,1,-1,0)$ is orthogonal to $r_1$, $r_2$
and $e_s$, $\nu=2$ exactly, while $L=32/h^2$. The left measure with masses
$\frac18,\frac38,\frac38,\frac18$ at $(s/h,u_1,u_2)=(0,0,0),(\frac12,1,0),
(\frac12,0,1),(1,1,1)$ and the right measure with masses $\frac12,\frac12$
at $(s/h,v)=(\frac14,0),(\frac34,1)$ have equal moments of $s/h$ through order
three ($\frac12,\frac5{16},\frac7{32}$) but not four ($\frac{11}{64}$ versus
$\frac{41}{256}$). Their value is $3/32$, a gap of $5/32$.

These bag moments even admit a global positive semidefinite completion.
With mean $(h/2,\frac12,\frac12,\frac12)$ and covariance
$\frac1{16}aa^{\T}+\frac3{16}bb^{\T}$, $a=(h,1,1,2)$,
$b=(0,1,-1,0)$, the completion satisfies every McCormick inequality for
$s\in[0,2h]$ and the control boxes. It violates the valid inequality
$u_1u_2+v-u_1v-u_2v=(1-v)u_1u_2+v(1-u_1)(1-u_2)\ge0$, a triangle inequality
of the Boolean quadric polytope \cite{Padberg1989}, with pseudo-expectation
$-1/8$. This inequality couples variables on both sides of the separator.
\end{remark}

